\begin{frame}{Introduction}
L'équation de Rayleigh-Stokes est une équation aux dérivées 
partielles utilisée en dynamique des fluides pour modéliser le 
comportement des 
fluides de second grade. Une forme canonique de l'équation pour 
une fonction de vitesse $u$ est
\begin{equation*}
    \partial_t u - \nu \Delta u - \alpha \Delta \big(\partial_t u\big) = f
\end{equation*}
%où
%\begin{itemize}
%   \item $\nu$ est la viscosité cinématique du fluide, responsable de la diffusion classique,
%   \item $\alpha$ est le paramètre de Rayleigh-Stokes, lié aux propriétés viscoélastiques du fluide,
%   \item $f$ est un terme de forçage,
%   \item $\Delta \big(\partial_t u\big)$ est un terme dit de \og dispersion régularisée \fg{} qui caractérise le comportement du fluide sous l'effet d'une accélération.
%\end{itemize}
\end{frame}